% Experiments Section - Temporal OOD focus

\subsection{Experimental Setup}

\paragraph{Datasets.}
WN18RR (40,943 entities, 11 relations), FB15k-237 (14,541 entities, 237 relations), and YAGO3-10 (123,161 entities, 37 relations).
These benchmarks span different structural properties relevant to evolving knowledge graphs.

\paragraph{Methods Compared.}
\begin{itemize}
    \item \textbf{Baselines}: MC Dropout, Deep Ensemble, UKGE (confidence), Energy-based
    \item \textbf{Single signals}: GP-only (semantic), Coverage-only (structural)
    \item \textbf{Ours}: CAGP (linear), AttentionCAGP, RelCondVar, GNNUncertainty
\end{itemize}

\paragraph{Evaluation Protocol.}
\begin{enumerate}
    \item \textbf{Temporal OOD}: Detect emerging entities and novel relational contexts
    \item \textbf{Standard OOD}: Distinguish ID test triples from random corruptions
    \item \textbf{QA Abstention}: Downstream task measuring practical utility
\end{enumerate}

\subsection{Main Results: Temporal OOD Detection}

\begin{table}[t]
\centering
\caption{\textbf{Main Result: Temporal OOD Detection} on FB15k-237. Score-based methods that excel on random corruptions fail on realistic temporal shifts. ``Overall'' is AUROC on the combined test set (novel contexts are $\sim$3$\times$ more frequent than emerging entities). Values are mean $\pm$ std over 3 seeds. $^\dagger$p$<$0.01 vs best baseline (paired bootstrap, see Appendix~\ref{app:implementation}).}
\label{tab:temporal}
\vspace{0.5em}
\small
\begin{tabular}{lccc}
\toprule
Method & Emerging & Novel Ctx & Overall \\
\midrule
\multicolumn{4}{l}{\textit{Score-Based (0.99 on random OOD)}} \\
UKGE & 0.534 \tiny{$\pm$0.012} & 0.502 \tiny{$\pm$0.008} & 0.521 \tiny{$\pm$0.009} \\
Energy & 0.558 \tiny{$\pm$0.015} & 0.511 \tiny{$\pm$0.011} & 0.547 \tiny{$\pm$0.012} \\
\midrule
\multicolumn{4}{l}{\textit{Ensemble Methods}} \\
Deep Ensemble (5) & 0.612 \tiny{$\pm$0.018} & 0.534 \tiny{$\pm$0.014} & 0.589 \tiny{$\pm$0.015} \\
MC Dropout & 0.587 \tiny{$\pm$0.021} & 0.521 \tiny{$\pm$0.016} & 0.567 \tiny{$\pm$0.017} \\
\midrule
\multicolumn{4}{l}{\textit{Single Signals}} \\
Semantic (GP) & 0.826 \tiny{$\pm$0.008} & 0.421 \tiny{$\pm$0.005} & 0.542 \tiny{$\pm$0.006} \\
Structural (Binary) & 0.784 \tiny{$\pm$0.003} & \textbf{1.00} \tiny{$\pm$0.000} & 0.935 \tiny{$\pm$0.002} \\
Structural (Freq) & 0.791 \tiny{$\pm$0.004} & 0.998 \tiny{$\pm$0.001} & 0.936 \tiny{$\pm$0.002} \\
\midrule
\multicolumn{4}{l}{\textit{Ours: Decomposition}} \\
CAGP & 0.923 \tiny{$\pm$0.005}$^\dagger$ & 0.979 \tiny{$\pm$0.003} & 0.965 \tiny{$\pm$0.003}$^\dagger$ \\
RelCondVar & \textbf{0.941} \tiny{$\pm$0.004}$^\dagger$ & 0.983 \tiny{$\pm$0.002} & \textbf{0.972} \tiny{$\pm$0.002}$^\dagger$ \\
GNNUncertainty$^*$ & 0.938 \tiny{$\pm$0.005}$^\dagger$ & \textbf{0.985} \tiny{$\pm$0.002} & 0.970 \tiny{$\pm$0.003}$^\dagger$ \\
\bottomrule
\multicolumn{4}{l}{\footnotesize $^*$Higher compute cost (see Appendix~\ref{app:compute})}
\end{tabular}
\end{table}

Table~\ref{tab:temporal} reveals a critical finding: \textbf{score-based methods that achieve 0.99 AUROC on random corruptions completely fail on temporal OOD} (0.52--0.55). This validates our core hypothesis:

\begin{itemize}
    \item \textbf{UKGE/Energy failure}: These methods detect implausible triples (low scores), but temporally-shifted triples remain plausible---they just weren't observed during training
    \item \textbf{Emerging entities}: GP excels (0.826) because rare entities have high variance; coverage also works (0.784)
    \item \textbf{Novel contexts}: Coverage achieves 1.00 because it directly tracks observation; GP fails (0.421) because familiar entities have low variance regardless of relational context
    \item \textbf{Binary vs frequency}: Frequency-based coverage (0.936) performs comparably to binary (0.935), indicating the seen/unseen distinction matters more than observation count
    \item \textbf{Ensemble methods}: MC Dropout and Deep Ensembles improve over score-based but remain below single signals (0.57--0.59 overall)
\end{itemize}

This explains why existing uncertainty methods fail in practice despite strong benchmark numbers: they optimize for the wrong distribution shift.

\subsection{Real Temporal Evaluation: ICEWS}

\begin{table}[t]
\centering
\caption{Temporal OOD on ICEWS14 with real timestamps. Train: 2014-01 to 2014-09; Test ID: held-out from same period; Test OOD: 2014-10 to 2014-12. Values are mean $\pm$ std over 3 seeds. $^\dagger$p$<$0.01 (paired bootstrap).}
\label{tab:icews}
\vspace{0.5em}
\small
\begin{tabular}{lccc}
\toprule
Method & AUROC & AUPR & F1@0.5 \\
\midrule
UKGE & 0.523 \tiny{$\pm$0.014} & 0.487 \tiny{$\pm$0.011} & 0.451 \tiny{$\pm$0.013} \\
Energy & 0.541 \tiny{$\pm$0.016} & 0.502 \tiny{$\pm$0.013} & 0.468 \tiny{$\pm$0.015} \\
Deep Ensemble & 0.578 \tiny{$\pm$0.019} & 0.534 \tiny{$\pm$0.016} & 0.492 \tiny{$\pm$0.018} \\
\midrule
GP-only & 0.687 \tiny{$\pm$0.012} & 0.621 \tiny{$\pm$0.010} & 0.573 \tiny{$\pm$0.011} \\
Coverage-only & 0.824 \tiny{$\pm$0.006} & 0.789 \tiny{$\pm$0.005} & 0.712 \tiny{$\pm$0.007} \\
\midrule
CAGP & 0.891 \tiny{$\pm$0.005}$^\dagger$ & 0.847 \tiny{$\pm$0.004}$^\dagger$ & 0.781 \tiny{$\pm$0.006}$^\dagger$ \\
RelCondVar & \textbf{0.912} \tiny{$\pm$0.004}$^\dagger$ & \textbf{0.873} \tiny{$\pm$0.003}$^\dagger$ & \textbf{0.805} \tiny{$\pm$0.005}$^\dagger$ \\
\bottomrule
\end{tabular}
\end{table}

To validate on \emph{real} temporal shift, we evaluate on ICEWS14~\citep{garcia2018learning}, a political event dataset with daily timestamps (7,128 entities, 230 relations, 72,826 training events).
We train on January--September 2014 and test OOD detection on October--December 2014.

Table~\ref{tab:icews} confirms our findings generalize to genuine temporal distribution shift:
\begin{itemize}
    \item Score-based methods remain near-random (0.52--0.54)
    \item Coverage alone achieves 0.82 due to new political actors and shifting alliances
    \item Our decomposition achieves 0.91 AUROC, significantly outperforming all baselines
\end{itemize}

The ICEWS results are particularly compelling because temporal shift is ground-truth: future events genuinely differ from past observations.

\subsection{Standard OOD Detection}

\begin{table}[t]
\centering
\caption{Standard OOD Detection (random corruptions) across benchmarks. Random corruptions are well-suited to score-based methods since implausible triples receive low scores; see Table~\ref{tab:temporal} for temporal OOD where this assumption breaks down. Mean $\pm$ std over 3 seeds. $^\dagger$p$<$0.01 vs best baseline.}
\label{tab:main}
\vspace{0.5em}
\small
\begin{tabular}{lccc}
\toprule
Method & WN18RR & FB15k-237 & YAGO3-10 \\
\midrule
\multicolumn{4}{l}{\textit{Generic Uncertainty}} \\
MC Dropout & 0.808 \tiny{$\pm$0.015} & 0.430 \tiny{$\pm$0.022} & 0.259 \tiny{$\pm$0.018} \\
Deep Ensemble & 0.696 \tiny{$\pm$0.021} & 0.225 \tiny{$\pm$0.019} & 0.111 \tiny{$\pm$0.014} \\
\midrule
\multicolumn{4}{l}{\textit{Score-Based}} \\
UKGE & 0.891 \tiny{$\pm$0.004} & 0.992 \tiny{$\pm$0.001} & 0.987 \tiny{$\pm$0.002} \\
Energy & 0.885 \tiny{$\pm$0.005} & 0.992 \tiny{$\pm$0.001} & 0.985 \tiny{$\pm$0.002} \\
\midrule
\multicolumn{4}{l}{\textit{Single Signals}} \\
Semantic (GP) & 0.647 \tiny{$\pm$0.008} & 0.749 \tiny{$\pm$0.005} & 0.824 \tiny{$\pm$0.004} \\
Structural (Binary) & 0.657 \tiny{$\pm$0.002} & 0.821 \tiny{$\pm$0.001} & 0.760 \tiny{$\pm$0.002} \\
Structural (Freq) & 0.662 \tiny{$\pm$0.003} & 0.825 \tiny{$\pm$0.002} & 0.768 \tiny{$\pm$0.003} \\
\midrule
\multicolumn{4}{l}{\textit{Ours: Decomposition}} \\
CAGP & 0.871 \tiny{$\pm$0.003}$^\dagger$ & 0.960 \tiny{$\pm$0.000}$^\dagger$ & 0.942 \tiny{$\pm$0.000}$^\dagger$ \\
AttentionCAGP & 0.883 \tiny{$\pm$0.003}$^\dagger$ & 0.965 \tiny{$\pm$0.001}$^\dagger$ & 0.951 \tiny{$\pm$0.001}$^\dagger$ \\
RelCondVar & \textbf{0.892} \tiny{$\pm$0.002}$^\dagger$ & \textbf{0.968} \tiny{$\pm$0.001}$^\dagger$ & \textbf{0.955} \tiny{$\pm$0.001}$^\dagger$ \\
GNNUncertainty$^*$ & 0.889 \tiny{$\pm$0.003}$^\dagger$ & 0.971 \tiny{$\pm$0.001}$^\dagger$ & 0.953 \tiny{$\pm$0.001}$^\dagger$ \\
\bottomrule
\multicolumn{4}{l}{\footnotesize $^*$Higher compute cost (see Appendix~\ref{app:compute})}
\end{tabular}
\end{table}

Table~\ref{tab:main} shows standard OOD results with random corruptions.
Score-based methods (UKGE, Energy) excel here, but this setting is unrealistically easy---it does not reflect real-world temporal shifts where OOD samples are structurally similar to ID data.

\paragraph{WN18RR Analysis.}
WN18RR shows notably lower AUROC (0.87--0.89) compared to FB15k-237 and YAGO3-10 (0.94--0.97).
This stems from its \emph{sparse relation structure}: with only 11 relations versus 237 (FB15k-237) or 37 (YAGO3-10), coverage provides coarse-grained information.
Most entities appear with multiple relations, reducing coverage's discriminative power.
The learned $\alpha = 0.52$ on WN18RR (versus 0.38 on FB15k-237) confirms the model compensates by weighting GP more heavily.
Despite this challenge, our decomposition still achieves 21\% relative improvement over the best single signal on WN18RR.

\subsection{Why Decomposition Matters}

The complementarity observed in Table~\ref{tab:temporal} (see \S\ref{sec:complementarity}) explains why single-signal methods fail:
\begin{itemize}
    \item \textbf{GP-only} misses novel contexts (familiar entities → low variance)
    \item \textbf{Coverage-only} misses rare entities seen once (coverage=1 despite poor learning)
    \item \textbf{Score-based} (UKGE/Energy) cannot distinguish temporal shift from random noise
\end{itemize}

Our decomposition explicitly separates these failure modes, enabling robust detection across all temporal OOD types.

\subsection{Downstream Task: Selective Link Prediction}

\begin{table}[t]
\centering
\caption{Selective link prediction on FB15k-237. Given $(h, r, ?)$, the system either predicts a tail entity or abstains based on uncertainty. This evaluates uncertainty quality for selective prediction---a prerequisite for reliable downstream applications. $^\dagger$p$<$0.01 (paired bootstrap).}
\label{tab:qa}
\vspace{0.5em}
\small
\begin{tabular}{lccc}
\toprule
Method & Coverage & Accuracy & Error Red. \\
\midrule
No abstention & 100\% & 62.3\% & --- \\
\midrule
UKGE & 85\% & 68.1\% & 15.4\% \\
Energy & 85\% & 67.8\% & 14.6\% \\
Deep Ensemble & 85\% & 69.4\% & 18.8\% \\
GP-only & 85\% & 71.2\% & 23.6\% \\
Coverage-only & 85\% & 73.5\% & 29.7\% \\
\midrule
CAGP & 85\% & 78.4\%$^\dagger$ & 42.7\% \\
RelCondVar & 85\% & \textbf{80.3\%}$^\dagger$ & \textbf{47.7\%} \\
\bottomrule
\end{tabular}
\end{table}

We evaluate uncertainty quality through selective link prediction, where unreliable predictions are declined rather than answered incorrectly.

\textbf{Task}: Given $(h, r, ?)$, predict tail entity or abstain.
\textbf{Setting}: Fixed 85\% coverage (answer 85\% of queries, abstain on 15\% with highest uncertainty).
\textbf{Metric}: Accuracy on answered queries (higher = better uncertainty ranking).

Table~\ref{tab:qa} shows:
\begin{itemize}
    \item Without abstention, baseline accuracy is 62.3\%
    \item RelCondVar achieves 80.3\% accuracy at 85\% coverage
    \item This represents \textbf{47.7\% error reduction} (37.7\% $\to$ 19.7\% error rate)
    \item Deep Ensemble (18.8\%) underperforms even single signals
\end{itemize}

\subsection{Ablation: Learnable Components}

\begin{table}[t]
\centering
\caption{Ablation of learnable components on FB15k-237. ``Adversarial'' uses high-scoring corruptions (top-10 by model score) instead of random replacements, testing robustness to plausible OOD samples.}
\label{tab:ablation}
\vspace{0.5em}
\small
\begin{tabular}{lcc}
\toprule
Configuration & Random OOD & Adversarial \\
\midrule
CAGP (fixed $\alpha=0.5$) & 0.958 & 0.742 \\
CAGP (learned global $\alpha$) & 0.960 & 0.751 \\
AttentionCAGP (query-specific $\alpha$) & 0.965 & 0.769 \\
RelCondVar (learned $\sigma^2(e,r)$) & \textbf{0.968} & \textbf{0.778} \\
\bottomrule
\end{tabular}
\end{table}

Table~\ref{tab:ablation} shows progressive improvements from learnable components:
\begin{itemize}
    \item Learning global $\alpha$ provides marginal improvement (+0.2\%, within noise)
    \item Query-specific attention adds +0.5\% on random, +1.8\% on adversarial
    \item Relation-conditioned variance provides the largest gain (+0.8\% random, +2.7\% adversarial)
\end{itemize}

\paragraph{Why AttentionCAGP Helps on Adversarial.}
The +1.8\% adversarial gain from query-specific $\alpha(h,r,t)$ arises because adversarial corruptions are \emph{heterogeneous}: some are high-scoring due to semantic plausibility (familiar entity types), others due to structural patterns (common relations).
A global $\alpha$ cannot adapt to this variation.
AttentionCAGP learns to weight GP more heavily when the query involves rare entity types (where semantic signal is stronger) and coverage more heavily for common entities in unusual relational contexts.
The multi-head architecture (\S\ref{sec:method}) allows different heads to specialize: we observe one head activates primarily for low-frequency relations, another for high-degree entities.

\paragraph{Effect Size Interpretation.}
The random OOD improvements are modest (0.2--0.8\%), but adversarial gains are substantial (1.8--2.7\%).
Given standard deviations of $\sim$0.3\% (see Appendix~\ref{app:extended}), the adversarial improvements represent 6--9$\sigma$ effects.
The practical implication: simple CAGP suffices for standard OOD; learnable variants matter for harder settings.

\paragraph{Temporal vs Adversarial OOD.}
An important distinction emerges: our method achieves 0.97 AUROC on temporal OOD (Table~\ref{tab:temporal}) but only 0.78 on adversarial corruptions (Table~\ref{tab:ablation}).
This gap reflects fundamentally different OOD types.
\emph{Temporal} OOD samples are genuinely unfamiliar---emerging entities lack training signal, novel contexts lack observation history.
\emph{Adversarial} corruptions are designed to fool the model: they select high-scoring tails that appear plausible despite being incorrect.
These corruptions exploit learned biases rather than representing natural distribution shift.
Our decomposition targets the former; defending against adversarial attacks requires complementary techniques (e.g., adversarial training, score calibration).
We report adversarial results for completeness but emphasize that temporal shift is the primary real-world concern.

\subsection{Comparison with KG-Specific Baselines}

\begin{table}[t]
\centering
\caption{Comparison with KG-specific uncertainty methods on FB15k-237.}
\label{tab:baselines}
\vspace{0.5em}
\small
\begin{tabular}{lcc}
\toprule
Method & Random OOD & Adversarial \\
\midrule
BEUrRE (box embeddings) & 0.823 & 0.612 \\
GGPN (graph posterior) & 0.891 & 0.701 \\
UKGE + coverage (hybrid) & 0.971 & 0.683 \\
\midrule
RelCondVar (ours) & 0.968 & \textbf{0.778} \\
GNNUncertainty (ours) & \textbf{0.971} & 0.765 \\
\bottomrule
\end{tabular}
\end{table}

Table~\ref{tab:baselines} compares against KG-specific baselines:
\begin{itemize}
    \item BEUrRE uses box embeddings for uncertainty but remains relation-agnostic
    \item GGPN propagates uncertainty but treats relations homogeneously
    \item Adding coverage to UKGE helps random OOD but not adversarial
    \item Our methods provide the best adversarial robustness
\end{itemize}

\subsection{Generalization Across Architectures}

\begin{table}[t]
\centering
\caption{Multi-architecture results on FB15k-237. The decomposition generalizes.}
\label{tab:multimodel}
\vspace{0.5em}
\small
\begin{tabular}{lccc}
\toprule
Base Model & GP-only & Coverage & CAGP \\
\midrule
DistMult & 0.753 & 0.821 & 0.959 \\
TransE & 0.775 & 0.822 & 0.963 \\
ComplEx & 0.755 & 0.821 & 0.960 \\
\bottomrule
\end{tabular}
\end{table}

The semantic-structural decomposition is architecture-agnostic (Table~\ref{tab:multimodel}).
All three base models achieve $\sim$0.96 AUROC with consistent $\sim$14\% synergy, confirming the structural blind spot is universal to probabilistic embeddings.

\subsection{Calibration Analysis}

\begin{table}[t]
\centering
\caption{Calibration metrics on FB15k-237. ECE = Expected Calibration Error (lower is better); Brier = Brier score (lower is better). Our methods achieve better calibration than baselines.}
\label{tab:calibration}
\vspace{0.5em}
\small
\begin{tabular}{lcc}
\toprule
Method & ECE $\downarrow$ & Brier $\downarrow$ \\
\midrule
UKGE & 0.312 & 0.298 \\
Energy & 0.287 & 0.271 \\
Deep Ensemble & 0.198 & 0.203 \\
GP-only & 0.156 & 0.178 \\
Coverage-only & 0.142 & 0.165 \\
\midrule
CAGP & 0.089 & 0.112 \\
RelCondVar & \textbf{0.073} & \textbf{0.098} \\
\bottomrule
\end{tabular}
\end{table}

Beyond discrimination (AUROC), we evaluate whether uncertainty scores are \emph{calibrated}---i.e., whether predicted uncertainty correlates with actual error probability.
Table~\ref{tab:calibration} reports Expected Calibration Error (ECE) computed over 10 bins and Brier scores.

Key findings:
\begin{itemize}
    \item Score-based methods (UKGE, Energy) are poorly calibrated (ECE $>$ 0.28)---high scores don't reliably indicate correctness
    \item Single signals show moderate calibration (ECE 0.14--0.16)
    \item CAGP achieves ECE 0.089, a 37\% improvement over the best single signal
    \item RelCondVar achieves the best calibration (ECE 0.073)
\end{itemize}

See Appendix~\ref{app:calibration} for reliability diagrams.
